\documentclass[11pt]{article}
\usepackage[margin=1in]{geometry}
\usepackage{amsmath,amssymb,amsthm,mathtools}
\usepackage{graphicx}
\usepackage{microtype}
\usepackage[hidelinks]{hyperref}

\newtheorem{theorem}{Theorem}
\newtheorem{corollary}{Corollary}
\newtheorem{lemma}{Lemma}
\newcommand{\N}{\mathbb N}
\newcommand{\R}{\mathbb R}

\title{Positive bidiagonal factorization for every\\finite-band oscillatory Toeplitz--Hessenberg matrix}
\author{Candidate full solution to the question in arXiv:2210.10728}
\date{12 August 2026}

\begin{document}
\maketitle

\noindent\textbf{Status.} Candidate full solution, pending expert review.

\medskip
\noindent\textbf{Source.} A.~Branquinho, A.~Foulquié-Moreno, and M.~Mañas,
\emph{Positive bidiagonal factorization of tetradiagonal Hessenberg
matrices}~\cite{source}, question on printed page~21.

\section{The question}
For a lower Hessenberg matrix with $r$ subdiagonals, a \emph{positive
bidiagonal factorization} (PBF) has the ordered form
\[
                         T=L_1L_2\cdots L_rU,                 \tag{1}\label{eq:pbf}
\]
where each $L_j$ is unit lower bidiagonal with strictly positive
subdiagonal, while $U$ is upper bidiagonal with unit superdiagonal and
strictly positive diagonal.  The source proves the Toeplitz case for $r=2$
and asks in its conclusion:

\begin{figure}[ht]
\centering
\includegraphics[width=\textwidth]{figures/open_problem_crop.png}
\caption{The exact question in the original paper, printed page~21.  This is
a full-width crop rendered from the archived source PDF.}
\end{figure}

\paragraph{Transcription.}
``Similar questions open for the constant Toeplitz matrix, do a oscillatory
pentadiagonal Toeplitz matrix always admit a PBF?''

The answer is yes.  In fact, the number of subdiagonals is arbitrary.

\section{Proof intuition}
The Edrei--Schoenberg classification turns total nonnegativity of a banded
Toeplitz matrix into a factorization of a polynomial:
\[
       1+e_1z+\cdots+e_{r+1}z^{r+1}
       =\prod_{j=1}^{r+1}(1+\beta_jz),\qquad \beta_j>0.
\]
Writing $S$ for the unilateral lower shift immediately factors the matrix as
one upper factor followed by $r$ lower factors.  This is exactly the wrong
order for \eqref{eq:pbf}.

The upper factor can nevertheless be moved rightward, one lower factor at a
time, without losing positivity.  After the first $m$ roots have been
absorbed, its $n$th diagonal entry is
\[
 q_n^{(m)}=\frac{h_{n+1}(\beta_1,\ldots,\beta_m)}
                  {h_n(\beta_1,\ldots,\beta_m)},
\]
where $h_n$ is the complete homogeneous symmetric polynomial.  The elementary
recurrence for $h_n$ is precisely the quotient--difference identity needed
to exchange an upper and a lower bidiagonal factor.  All $h_n$ are positive,
so every new factor is positive.  Iterating gives the requested order.

\section{Statement and proof}
Index rows and columns by $\N_0$.  Let $S$ be the unilateral lower shift,
$S_{n,n-1}=1$ for $n\ge1$; thus $S^*S=I$.  For $r\ge1$ consider
\[
 T=S^*+e_1I+e_2S+\cdots+e_{r+1}S^r.             \tag{2}\label{eq:T}
\]
This is a Toeplitz--Hessenberg matrix with $r+2$ nonzero diagonals.

\begin{theorem}\label{thm:all}
For every $r\ge1$, an oscillatory Toeplitz--Hessenberg matrix
\eqref{eq:T} admits a PBF \eqref{eq:pbf}.  More explicitly, if
\[
 1+e_1z+\cdots+e_{r+1}z^{r+1}
       =\prod_{j=1}^{r+1}(1+\beta_jz),\qquad \beta_j>0,       \tag{3}\label{eq:edrei}
\]
then the entries of all factors are given by positive quotients of complete
homogeneous polynomials, as in \eqref{eq:lambda} below.
\end{theorem}

\begin{proof}
The Edrei--Schoenberg theorem~\cite{edrei,schoenberg}, applied exactly as in
the source's tetradiagonal Toeplitz argument, says that oscillation of
\eqref{eq:T} yields \eqref{eq:edrei}; positivity of the outermost
subdiagonal makes every $\beta_j$ strictly positive.

For $m\in\{1,\ldots,r+1\}$ and $n\ge0$, set
\[
 H_n^{(m)}:=h_n(\beta_1,\ldots,\beta_m),
 \qquad q_n^{(m)}:=\frac{H_{n+1}^{(m)}}{H_n^{(m)}}.
\]
Here $H_0^{(m)}=1$, and every $H_n^{(m)}>0$.  Let $U^{(m)}$ be the upper
bidiagonal matrix with diagonal $(q_n^{(m)})_{n\ge0}$ and unit
superdiagonal.  In particular,
$U^{(1)}=S^*+\beta_1I$.  Write $L(\beta)=I+\beta S$.

For $m\ge2$, define a unit lower bidiagonal matrix $B_{m-1}$ by
\[
 (B_{m-1})_{n,n-1}=\lambda_n^{(m)}
 :=\frac{\beta_m q_n^{(m-1)}}{q_{n-1}^{(m)}}>0,
 \qquad n\ge1.                                      \tag{4}\label{eq:lambda}
\]
We claim the local exchange identity
\[
                 U^{(m-1)}L(\beta_m)=B_{m-1}U^{(m)}. \tag{5}\label{eq:exchange}
\]
Both sides have unit superdiagonal and no entries below the first
subdiagonal.  Their $(n,n-1)$ entries agree by \eqref{eq:lambda}.  At $n=0$,
the diagonal equality is
$q_0^{(m)}=q_0^{(m-1)}+\beta_m$.  For $n\ge1$, use
\[
 H_n^{(m)}-\beta_mH_{n-1}^{(m)}=H_n^{(m-1)},
 \qquad
 H_{n+1}^{(m)}=H_{n+1}^{(m-1)}+\beta_mH_n^{(m)}.    \tag{6}\label{eq:hrec}
\]
Indeed,
\begin{align*}
 q_n^{(m-1)}+\beta_m-\lambda_n^{(m)}
 &=\frac{H_{n+1}^{(m-1)}}{H_n^{(m-1)}}
   \left(1-\frac{\beta_mH_{n-1}^{(m)}}{H_n^{(m)}}\right)+\beta_m\\
 &=\frac{H_{n+1}^{(m-1)}}{H_n^{(m)}}+\beta_m
  =\frac{H_{n+1}^{(m)}}{H_n^{(m)}}=q_n^{(m)}.
\end{align*}
This proves \eqref{eq:exchange}, including strict positivity of its new
lower factor.

Finally, $S^*S^k=S^{k-1}$ for $k\ge1$, so expansion of \eqref{eq:edrei}
gives the exact initial factorization
\[
 T=U^{(1)}L(\beta_2)L(\beta_3)\cdots L(\beta_{r+1}). \tag{7}\label{eq:wrong}
\]
Apply \eqref{eq:exchange} successively for $m=2,3,\ldots,r+1$:
\[
 T=B_1U^{(2)}L(\beta_3)\cdots L(\beta_{r+1})
   =\cdots=B_1B_2\cdots B_rU^{(r+1)}.
\]
Every subdiagonal entry of every $B_j$ and every diagonal entry of
$U^{(r+1)}$ is strictly positive.  This is the required PBF.
\end{proof}

\begin{corollary}[answer to the source question]
Every oscillatory pentadiagonal Toeplitz--Hessenberg matrix admits a positive
bidiagonal factorization with three unit lower bidiagonal factors followed by
one positive upper bidiagonal factor.
\end{corollary}

\begin{proof}
Take $r=3$ in Theorem~\ref{thm:all}.
\end{proof}

\section{Verification, scope, and novelty}
The exact-arithmetic script \texttt{code/verify\_pbf.py} constructs the
factors from \eqref{eq:lambda} and compares their product entrywise with
twelve-by-twelve leading Toeplitz sections.  It checks five cases from one
through five subdiagonals, including repeated roots.  All checks pass.  This
is a normalization check only; equations \eqref{eq:hrec}--\eqref{eq:wrong}
constitute the proof.

The result concerns the source's one-superdiagonal Toeplitz--Hessenberg
setting.  It does not answer the source's broader final question about
matrices having several superdiagonals as well as several subdiagonals, nor
does it claim that arbitrary non-Toeplitz oscillatory Hessenberg matrices have
a PBF.

\paragraph{Novelty check.}
On 12 August 2026, bounded arXiv and run-index searches used the exact title
and source question together with ``pentadiagonal Toeplitz PBF'', ``banded
Toeplitz bidiagonal factorization'', ``quotient--difference'', ``complete
homogeneous'', and ``branched continued fractions''.  Later
arXiv:2412.03694 and arXiv:2603.21345 develop broad MOP/branched
factorizations, but source inspection and exact searches found neither the
Toeplitz theorem nor formula \eqref{eq:lambda}.  Novelty confidence is
moderate rather than definitive.

\paragraph{Human-review recommendation.}
Prioritize review.  The main checks are the orientation of the source's PBF
convention and the general finite-polynomial Edrei--Schoenberg invocation;
the quotient--difference calculation itself is elementary and exact.

\begin{thebibliography}{9}
\bibitem{source}
A.~Branquinho, A.~Foulquié-Moreno, and M.~Mañas,
\emph{Positive bidiagonal factorization of tetradiagonal Hessenberg
matrices}, arXiv:2210.10728 (2022).

\bibitem{edrei}
A.~Edrei, \emph{Proof of a conjecture of Schoenberg on the generating
function of a totally positive sequence}, Canad. J. Math. \textbf{5} (1953),
86--94.

\bibitem{schoenberg}
I.~J. Schoenberg, \emph{Some analytical aspects of the problem of smoothing},
in \emph{Studies and Essays Presented to R. Courant on His 60th Anniversary},
Interscience, 1948, 351--370.
\end{thebibliography}

\end{document}
